% ============================================================
% SNIPPET FINAL PARA CAPÍTULO 5 — Validación del gateway
% Datos reales: 17 visitas Mayo 11-27, 33 sesiones con correlación
% ============================================================

\subsection{Validación de laboratorio y verificación en campo}
\label{subsec:validacion_gateway_campo}

La descripción de la arquitectura en las secciones anteriores cubre el diseño
declarado; esta sección documenta la verificación empírica de que la pasarela
se comporta conforme a ese diseño en condiciones reales de campo.

La suite se dividió en dos planos complementarios. Las pruebas de
normalización (esquema \acs{JSON}, idempotencia y consistencia de
$\eta$) se ejecutaron mediante dry-run en Windows sobre las 33
sesiones de ordeño completo del período 11--28 de mayo de 2026.
Las pruebas de canal seguro y resiliencia operativa (\acs{TLS},
recuperación del spool ante caída del broker, recepción confirmada
por suscriptor externo) se ejecutaron sobre la Raspberry Pi en 17
visitas del mismo período (11--27 de mayo de 2026).

\begin{table}[H]
\centering
\caption{Resultados de la suite de validación de la pasarela perimetral
         (período post-intervención, 17 visitas, 33 sesiones, 11--28 de mayo de 2026).}
\label{tab:gateway_validation_suite}
\footnotesize
\renewcommand{\arraystretch}{1.22}
\begin{tabularx}{\linewidth}{|>{\raggedright\arraybackslash}p{3.4cm}|>{\raggedright\arraybackslash}X|>{\centering\arraybackslash}p{2.2cm}|>{\raggedright\arraybackslash}p{3.6cm}|}
\hline
\rowcolor{gray!20}
\TableHeaderFirst{Prueba} & \TableHeader{Alcance / plataforma} & \TableHeader{Resultado} & \TableHeader{Observación} \\ \hline

Esquema \acs{JSON} &
Estructura y tipos de campo; dry-run Windows por sesión &
\textbf{33/33} \acs{PASS} &
Verificado sobre cada sesión con correlación efectiva. \\ \hline

Idempotencia &
Reproducibilidad \acs{MD5} entre corridas aisladas; dry-run Windows &
\textbf{33/33} \acs{PASS} &
Directorios aislados por sesión; coincidencia exacta en re-ejecución. \\ \hline

Consistencia $\eta$ &
$\eta_{\text{gateway}} = \eta_{\text{motor}}$; dry-run Windows &
\textbf{33/33} \acs{PASS} &
$\bar{\eta}_{\text{post}} = 26{,}25\%$ (rango: $8{,}3\%$--$100{,}0\%$). \\ \hline

\acs{TLS} Handshake &
Cifrado del canal; verificado en Pi (17 visitas) &
\textbf{17/17} \acs{PASS} &
\acs{TLS} 1.2 rechazado; \acs{TLS} 1.3 aceptado con certificado de cliente. \\ \hline

Resiliencia del spool &
Continuidad ante caída del broker; verificado en Pi (17 visitas) &
\textbf{17/17} \acs{PASS} &
Lote completo drenado tras recuperación; spool residual $= 0$. \\ \hline

Suscripción \acs{MQTT} &
Recepción por suscriptor externo; verificado en Pi (17 visitas) &
\textbf{17/17} \acs{PASS} &
Mensajes verificados con \texttt{mosquitto\_sub} post-recuperación. \\ \hline

\end{tabularx}
\end{table}

La prueba de resiliencia se estructuró en cinco pasos secuenciales:
(i)~detener el broker \texttt{mosquitto}, (ii)~ejecutar el gateway
para que espole el lote en disco, (iii)~reiniciar el broker,
(iv)~activar un suscriptor externo e invocar el drenado, y
(v)~confirmar que el suscriptor recibe exactamente los mensajes
esperados y que el spool queda vacío. El
Listado~\ref{lst:resiliencia_salida} muestra la salida
representativa sobre una sesión de cuatro mensajes
(Visita~25, 25/05/2026).

\begin{lstlisting}[
  language=bash,
  caption={Salida del test de resiliencia en \texttt{gateway-esmeralda}
           (29/05/2026, sesión Visita~25).},
  label=lst:resiliencia_salida,
  basicstyle=\ttfamily\footnotesize,
  breaklines=true,
  frame=single
]
SPOOL_FILES=1 SPOOL_LINES=4
drained=4 failed=0
RECEIVED=4 SPOOL_AFTER=0
\end{lstlisting}

Durante la verificación de resiliencia se identificó que la entrega
de mensajes con \acs{QoS}$\geq 1$ requería que el bucle de red del
cliente \texttt{paho-mqtt} permaneciera activo hasta la confirmación
del broker. Se corrigió \texttt{publish\_batch} añadiendo
\texttt{loop\_start()} antes de publicar y \texttt{loop\_stop()}
tras confirmar la entrega, garantizando que cada lote se transmite
completamente antes de cerrar la conexión. Esta corrección no
modifica el comportamiento del spool ni el protocolo de reintentos.

Con estas pruebas se cerró la verificación del componente de
transporte. El contraste estadístico de $\eta$ pre/post intervención
se presenta en el Capítulo~\ref{chapter:chapter06}.
